\section{Discusión crítica}
\subsection{¿Por qué no se alcanza el 80\,\% de goles?}
El éxito de las habilidades aisladas no garantiza su composición: definición
alcanza 86,0\,\% desde la zona de tiro, mientras el mejor ataque completo obtiene
65,0\,\% (IC 95\,\%: 58,2--71,3\,\%), por debajo de la meta del 80\,\%.
Entrenar definición dentro de la jerarquía eleva el resultado del sistema con
traspaso anticipado de 45,0\,\% a 65,0\,\%. Esto respalda la relevancia de adaptar
la habilidad a los estados que recibe, aunque el entrenamiento adicional impide
atribuir toda la mejora únicamente al cambio de distribución.

\subsection{¿Por qué el ajuste con RL no siempre mejora la imitación?}
No hay un algoritmo superior en todas las tareas. DQN aprende persecución con
menos interacciones que PPO y obtiene 83,0\,\% en definición tras imitación,
pero queda en 70,5\,\% en conducción, frente a 100\,\% de A2C y 98,0\,\% de PPO.
El ajuste del sistema completo tampoco supera al coordinador DAgger humano:
26,5\,\% con PPO frente a 30,5\,\%. Estos resultados no demuestran sobreajuste
del replay ni identifican por sí solos la causa del deterioro.

\subsection{¿Qué papel tienen la recompensa y el experto?}
La recompensa y las demostraciones también condicionan el aprendizaje.
En conducción, el shaping ingenuo eleva la recompensa recibida, pero alcanza
92,7\,\% de éxito frente a 99,2\,\% con PBRS. DAgger humano mejora BC de
13,0\,\% a 53,5\,\%, mientras el experto programado alcanza al menos 99,5\,\%.
La mejora depende de la tarea y del experto; no permite atribuir los fallos
humanos a tiempos de reacción o falta de datos sin un análisis adicional.

\subsection{¿Qué limita las conclusiones y cómo mejorar la evaluación?}
La evidencia más sólida entre semillas corresponde a la escalera de persecución;
las curvas de ajuste fino y las variantes jerárquicas no acreditan esa misma
estabilidad. Además, se usa percepción privilegiada, rivales heurísticos y una
condición estandarizada de prueba. El costo previo de imitación queda fuera del
presupuesto de ajuste con RL. El siguiente paso es repetir las mejores variantes
con varias semillas, igualar presupuestos y variar posición inicial y dificultad
del rival para medir generalización y separar el efecto de cada modificación.
